\section{总结与延伸}

\subsection{全文核心脉络}

本文描绘了 LLM 发展的三个阶段：

\begin{enumerate}
    \item \textbf{Pretraining Scaling 时代}：通过扩大预训练规模提升模型能力
    \item \textbf{Reasoning Thinking 时代}（o1/R1）：通过 RL 后训练让模型``思考更久''，核心依赖强反馈信号和基础设施
    \item \textbf{Agentic Thinking 时代}（当前转型）：从``思考更久''到``为行动而思考''，核心对象变为 model + environment 系统
\end{enumerate}

\begin{importantbox}{竞争优势的转移}
\textit{In the reasoning era, the edge came from better RL algorithms, stronger feedback signals, and more scalable training pipelines. In the agentic era, the edge will come from better environments, tighter train-serve integration, stronger harness engineering, and the ability to close the loop between a model's decisions and the consequences those decisions produce.}

在 reasoning 时代，竞争优势来自更好的 RL 算法、更强的反馈信号和更可扩展的训练流水线。在 agentic 时代，竞争优势将来自更好的环境、更紧密的训练-服务集成、更强的 harness engineering，以及闭合模型决策与其后果之间循环的能力。
\end{importantbox}

\subsection{关键启示}

\begin{itemize}
    \item ``好的思考''的定义已经改变：不再是最长或最可见的轨迹，而是在真实世界约束下维持有效行动的最有用的轨迹
    \item 研究工件的重心在转移：从模型架构和训练数据，扩展到环境设计、rollout 基础设施、评估器鲁棒性和多 agent 协调接口
    \item Thinking 与 Instruct 的合并仍是开放问题，``有机合并''（推理努力的连续谱）比``机械合并''（两种模式的拼接）更有价值
    \item 环境质量将取代数据多样性成为新时代的核心资产
\end{itemize}

\subsection{拓展阅读}

\begin{itemize}
    \item Qwen3 技术博客：混合思考模式与四阶段后训练流水线
    \item OpenAI o1 系统卡与技术报告
    \item DeepSeek-R1 技术报告
    \item Anthropic Claude 3.7 Sonnet / Claude 4 发布说明
    \item DeepSeek V3.1 Think \& Non-Think 混合推理
\end{itemize}

%% --- Fin del contenido --- %%

\end{document}
